\documentclass[aip,preprint,amsmath,amssymb]{revtex4-2}

\usepackage{graphicx}
\usepackage{booktabs}
\usepackage{mathptmx}
\graphicspath{{figures/}{paper_figures/}}

\begin{document}

\title[Death of a Chimera]{Death of a Chimera: Breath-Synchronized Collapse of a Post-Homoclinic Transient in a Two-Population Oscillator System}

\author{Avery Karlin}
\thanks{ORCID: 0000-0003-3848-6782}
\email{ak5232@columbia.edu}
\affiliation{Independent Researcher}

\date{\today}

\begin{abstract}
We report a long-duration empirical study of chimera collapse in the two-population
Sakaguchi--Kuramoto system with strictly identical natural frequencies, on the
oscillator engine of a deployed additive synthesizer at small population sizes
($N = 4$--$64$ per population). Its operational collapse detector --- and the
first-passage criterion standard for chimera lifetimes --- measure the wrong
observable: up to 98\% of sustained near-synchrony crossings recover as the chimera
re-forms, and 69--97\% of detector firings are self-healing \emph{grazes}.
Under an absorption-grade label, true lifetimes at the operating corner ($A = 0.5$,
$\beta = 0.05$) double, become $N$-independent ($\tau_{\mathrm{abs}} \approx 139$\,s
over a $16\times$ range of $N$), and age (Weibull shape $2.3$--$3.0$). Deaths are
breath-synchronized (pooled Rayleigh $p = 2.2\times10^{-6}$) and ratchet: per-cycle
envelope maxima grow monotonically to capture. The reduced two-population model of
Abrams \emph{et al.}, validated against its published bifurcations with no fitted
parameters, explains the inversion: the operating corner lies beyond the homoclinic
($A_{hc}(0.05) = 0.4096 < 0.5$), so the deployed ``chimera'' is a transient
spiraling out along a destroyed limit cycle's ghost, while the nominally unreliable
corner ($A = 0.2$) holds the only true attractor ($r^* = 0.678$; 51--78\%
of seeds never absorb). It matches the breath period (25.0 vs.\ 21--25\,s) and
spiral rates and predicts individual lifetimes beyond any static feature, yet
captures $3.2\times$ faster than the finite-$N$ system dies. Finite-size
fluctuations \emph{prolong} the transient --- an inversion of the
noise-kills-metastability intuition; a direct test rules out additive collective
noise, leaving the mechanism open.
\end{abstract}

\maketitle


\begin{quotation}
\noindent Chimera states, in which a population of identical oscillators
spontaneously splits into one synchronized group and one disordered group, are
usually studied as a question in pure mathematics. Here the same phenomenon appears
inside a working music synthesizer, where two coupled oscillator banks shape the
sound and the software periodically judges the pattern ``dead'' and restarts it.
Asking why it dies so often reveals that the restart trigger was watching the wrong
quantity: most apparent deaths are momentary lapses that heal on their own. Once the
genuine deaths are isolated, they prove regular, rhythmic, and no more frequent in
large populations than small ones --- and a simple reduced model explains where and
why they occur, while underpredicting how long the pattern survives, a discrepancy
the paper leaves as an open problem.
\end{quotation}

\section{Introduction}\label{sec:intro}

Chimera states --- the coexistence of synchronized and desynchronized populations in
systems of identical coupled oscillators --- have been a central object of nonlinear
dynamics since their discovery by Kuramoto and Battogtokh~\cite{kb2002} and naming by
Abrams and Strogatz~\cite{as2004}. The two-population model of Abrams, Mirollo,
Strogatz, and Wiley~\cite{abrams2008} made the phenomenon solvable: two identical
populations with stronger intra- than inter-population coupling support stable
chimeras, breathing chimeras born in a Hopf bifurcation, and a homoclinic boundary
beyond which no chimera attractor survives. At finite size the picture is enriched by
transience: on nonlocally coupled rings, chimeras are chaotic transients whose
lifetimes grow with oscillator count~\cite{wo2011}, and the finite-$N$ fate of
chimeras has remained a recurring theme of the literature
since~\cite{pa2015,om2018}, including for two small populations of the kind
studied here~\cite{panaggio2016}.

This paper studies chimera collapse in an unusual arena: the oscillator engine of a
deployed real-time additive synthesizer, in which a two-population
Sakaguchi--Kuramoto bank with \emph{strictly identical} natural frequencies drives
the partial amplitudes of a musical voice, and a supervisor process re-seeds the
chimera whenever an engineering detector declares it dead. Coupled-oscillator
networks have an established history as sound-synthesis drivers~\cite{lem2019,lemsmc};
here the synthesis application matters to the dynamics only in that it fixes the
regime of interest --- small populations ($N = 4$--$64$ per side), long unattended
runtimes, and a specific operating corner chosen, as it turns out, for the wrong
reason. The deployment prompts a question the theory literature rarely needs to ask
operationally: \emph{what exactly is the death of a chimera, and what does a detector
of it actually detect?}

We answer with a chain of measurements, each gating the next, on a bit-reproducible
integrator (2{,}100 baseline runs plus targeted re-campaigns; every figure
regenerable from committed configurations). The contributions are:

\begin{enumerate}
\item \textbf{Graze versus absorption.} The operational collapse criterion ---
sustained near-synchrony of the incoherent population --- overwhelmingly fires on
events the chimera survives: up to 98\% of criterion-satisfying crossings recover,
and 69--97\% of the deployed detector's firings are such grazes. We introduce a
two-timescale, absorption-grade labeling and show first-passage ``lifetimes'' are
roughly half the true absorption times (Sec.~\ref{sec:graze}).
\item \textbf{An inverted stability picture.} Under absorption labeling, lifetimes at
the operating corner are $N$-independent ($\approx 139$\,s, $N = 4$--$64$, every seed
absorbing), while at the nominally unreliable high-morph corner the majority of seeds
\emph{never} absorb: the operating corner has no chimera attractor at all, and the
discarded corner holds the only persistent chimera in the system
(Sec.~\ref{sec:stats}).
\item \textbf{The death mechanism.} Collapse is a breath-synchronized ratchet:
collapse attempts and true deaths lock to a phase of the breathing cycle, successive
per-cycle envelope maxima grow monotonically (spiraling out exponentially toward the
synchronization boundary), and hazard rises with cycle index. Transverse escape from
the Poisson submanifold is ruled out by a validated observable and a null
interventional test; lifetime is instead organized by the collective initial
condition, dominantly the inter-population phase offset (Sec.~\ref{sec:mechanism}).
\item \textbf{Closure by the reduced model --- and a clean open problem.} The reduced
equations of Ref.~\cite{abrams2008}, validated against that paper's own published
results and applied with a literal 1:1 timescale and zero fitted parameters, place
the operating corner beyond the homoclinic ($A_{hc}(\beta{=}0.05) = 0.4096$),
reproduce the breath period and per-cycle spiral rates, identify the never-absorbers'
level with the stable chimera's $r^*$, and predict individual run lifetimes beyond
any static feature of the seed. They also fail in one instructive way: mean-field
capture is $3.2\times$ faster than finite-$N$ death, and the flat
$\tau_{\mathrm{abs}}(N)$ is absent from the $N$-free collective flow --- finite-size
fluctuations \emph{extend} the transient (Sec.~\ref{sec:reduced},
Sec.~\ref{sec:discussion}).
\end{enumerate}

\section{System}\label{sec:system}

The voice integrates two populations $\sigma \in \{1, 2\}$ of $N$ phase oscillators
each,
\begin{equation}
\dot\theta_i^\sigma \;=\; \omega \;+\; \sum_{\sigma'} \frac{K_{\sigma\sigma'}}{N}
\sum_{j=1}^{N} \sin\!\big(\theta_j^{\sigma'} - \theta_i^{\sigma} - \alpha\big),
\label{eq:model}
\end{equation}
with intra-population coupling $K_{\sigma\sigma} = \mu = (1+A)/2$, inter-population
coupling $K_{\sigma\sigma'} = \nu = (1-A)/2$ ($\mu + \nu = 1$, $A = \mu - \nu$), and
phase lag $\beta = \pi/2 - \alpha$~\cite{kuramoto1984,sk1986}, matching the
conventions of Ref.~\cite{abrams2008}. A code audit confirms the natural frequencies are identical
\emph{exactly} ($\delta\omega = 0$: no per-oscillator frequency term, detuning
explicitly zeroed on this path), so the dynamics is deterministic given the initial
condition and collapse cannot be attributed to frequency disorder; we work in the
$\omega = 0$ rotating frame. Integration is fourth-order Runge--Kutta with
$dt = 0.05$ at a documented one-model-unit-per-second clock, so all times below are
seconds with no conversion; halving and quartering $dt$ leaves lifetime
distributions statistically indistinguishable (two-sample KS $p = 1.0$). All runs are
bit-reproducible from logged (seed, parameters, $dt$); every re-analysis below begins
by reproducing the upstream labels exactly.

In the application, the incoherent population's per-oscillator coherences modulate
the gains of an additive partial bank, $g_i' = g_i\,[\,1 + d\,a\,(c_i - \tfrac12)\,]$
with $c_i = \tfrac12(1 + \cos(\theta_i - \psi))$, turning the chimera's collective
dynamics into a slowly self-morphing spectrum; the synthesis layer plays no further
role in the dynamics studied here. Runs are initialized with the canonical chimera
seed (one population tightly clustered, the other spread), since spontaneous chimera
formation from random phases is negligible in this regime. The deployed supervisor
declares the state alive while $\max_\sigma R_\sigma > 0.9$ and
$\min_\sigma R_\sigma < 0.85$, fires after 2\,s of sustained violation, and re-seeds
the incoherent population. The baseline campaign comprises 200 seeds per population
size at the operating corner ($A = 0.5$) and 100 at the high-morph corner
($A = 0.2$), $N \in \{4, 8, 16, 24, 32, 48, 64\}$, $\beta = 0.05$ throughout,
integrated unsupervised to $t_{\max} = 2000$\,s.

\section{Graze versus absorption}\label{sec:graze}

What counts as the death of a chimera? In a deployed system the answer is
operational: the supervisor's aliveness condition fails continuously for 2\,s, and
our initial lifetime campaign used the closely related first-passage criterion that
$\min(R_1, R_2)$ exceed a threshold $\theta = 0.85$ sustained for $W = 5$\,s. Both
definitions encode the same assumption: that a sustained excursion of the incoherent
population to near-synchrony \emph{is} the death of the chimera. This section shows
that assumption is wrong, and replaces it.

\subsection{Threshold crossings overwhelmingly recover}

Extending the recorded trajectories past their first criterion-satisfying crossing,
rather than terminating there, reveals that the crossing is usually not absorbing. At
the operating corner and across $N = 8$--$64$, up to 98\% of first crossings ---
excursions exceeding $\theta$ for longer than the full 5\,s sustain window --- are
followed by recovery: the incoherent population desynchronizes again, with
post-crossing minima of $\min(R_1, R_2)$ clustering at 0.2--0.45, a fully re-formed
chimera. About two-thirds of these recovered trajectories are later absorbed for
good; the remainder churn through further graze--recover cycles. The first sustained
crossing is therefore a \emph{graze} of the synchronization boundary, not a death,
and first-passage lifetimes computed from it measure time-to-first-long-graze.

Figure~\ref{fig:trace} shows a representative trajectory ($N = 16$, $A = 0.5$; seed
pinned in the released configuration). The run grazes twice --- each time exceeding
$\theta$ for more than 5\,s before the chimera re-forms --- and is absorbed on its
third approach at $t_{\mathrm{abs}} = 221.9$\,s, having first ``collapsed'' by the
engineering criterion at $t_{\mathrm{graze}} = 83.2$\,s. The shipped supervisor's
detector fires \emph{six times} on this trajectory before the single event that is
irreversible: six firings, two reforms, one death.

\begin{figure}[t]
\centering
\includegraphics[width=0.95\linewidth]{fig1}
\caption{A graze is not a death. $\min(R_1, R_2)$ for a single run ($N = 16$,
$A = 0.5$, $\beta = 0.05$; breath period $T_b = 20.5$\,s). The trajectory crosses
$\theta = 0.85$ (upper line) in sustained excursions twice, recovering below the
0.80 recovery threshold (lower line) each time as the chimera re-forms, before the
absorbing crossing at $t_{\mathrm{abs}} = 221.9$\,s. The engineering first-passage
label $t_{\mathrm{graze}} = 83.2$\,s and the six firings of the deployed 2-s detector
are marked.}
\label{fig:trace}
\end{figure}

\subsection{An absorption-grade label}

We therefore assign every run two timescales from a single trace. The first,
$t_{\mathrm{graze}}$, is the unmodified first-passage label, retained for
comparability. The second, $t_{\mathrm{abs}}$, is the time of the first crossing
\emph{not} followed by recovery --- $\min(R_1, R_2) < 0.80$ sustained at least
5\,s --- within a verification horizon $T_v = 120$\,s; runs still alternating
between grazes and recoveries at the integration horizon are right-censored.

The absorption label is robust to its own knobs where it matters. Sweeping
$T_v \in \{60, 120, 240\}$\,s against recovery thresholds $\{0.75, 0.80\}$ moves
$\hat\tau_{\mathrm{abs}}$ by 22--27\% at $A = 0.5$, monotonically in $T_v$ and
negligibly in the recovery threshold --- comparable to or better than the
$\sim$52\% $(\theta, W)$-sensitivity of the graze criterion itself. At $A = 0.2$
the apparent sensitivity is dominated by censoring rather than the labeling
parameters, for a reason Sec.~\ref{sec:stats} shows is dynamical: most $A = 0.2$
trajectories never absorb at all.

The quantitative consequence at the operating corner is a factor of two:
$\hat\tau_{\mathrm{graze}} \approx 63$--$67$\,s against
$\hat\tau_{\mathrm{abs}} \approx 139$\,s (Fig.~\ref{fig:tau}). First-passage chimera
lifetimes computed with engineering-style criteria systematically measure the wrong,
and roughly halved, quantity.

\subsection{The deployed detector is a graze detector}

Replaying the shipped detector over the same unsupervised trajectories and
classifying each firing by whether the underlying event self-recovers gives
Table~\ref{tab:overtrigger}: 69--79\% of firings at $A = 0.5$, and 95--97\% at
$A = 0.2$, occur on grazes the chimera survives on its own. The deployed system has
been re-seeding a state that was mostly not dead --- harmless in practice, since
re-seeding lands near where the trajectory already was, but diagnostic of the
conceptual gap this section closes: engineering collapse detection and dynamical
absorption are different observables. All statistics below are computed on
$t_{\mathrm{abs}}$ unless noted.

\begin{table}[t]
\centering
\caption{Supervisor over-trigger: fraction of the deployed 2-s detector's firings,
replayed over unsupervised traces, occurring on self-recovering grazes.}
\label{tab:overtrigger}
\begin{tabular}{ccccc}
\toprule
$A$ & $N$ & firings/run & over-trigger & median wasted (s) \\
\midrule
0.5 & 8  & 7.2  & 79\% & 4.4 \\
0.5 & 16 & 5.2  & 76\% & 2.4 \\
0.5 & 32 & 4.1  & 73\% & 3.9 \\
0.5 & 64 & 3.3  & 69\% & 5.3 \\
0.2 & 8  & 12.3 & 95\% & 2.0 \\
0.2 & 16 & 9.8  & 97\% & 1.8 \\
0.2 & 32 & 7.9  & 96\% & 1.7 \\
0.2 & 64 & 4.7  & 95\% & 1.8 \\
\bottomrule
\end{tabular}
\end{table}

\begin{figure}[t]
\centering
\includegraphics[width=0.85\linewidth]{fig2}
\caption{Graze versus absorption lifetimes. Exponential-MLE
$\hat\tau_{\mathrm{graze}}$ and $\hat\tau_{\mathrm{abs}}$ versus $N$ at both
coupling disparities. At $A = 0.5$ absorption lifetimes are $\approx 2.2\times$ the
first-passage values and flat in $N$; at $A = 0.2$ no absorption lifetime is
defined for the majority of seeds (Sec.~\ref{sec:stats}).}
\label{fig:tau}
\end{figure}

\section{Absorption statistics at finite $N$}\label{sec:stats}

\subsection{A flat plateau at the operating corner}

At $A = 0.5$, every one of the 1{,}400 seeds absorbs within the horizon (0\%
censored at every $N$), and the absorption lifetime is $N$-independent:
$\hat\tau_{\mathrm{abs}} = 139, 152, 143, 131, 136, 136, 135$\,s for
$N = 4$--$64$, a ratio $\tau(64)/\tau(4) = 0.97$ over a sixteen-fold range of
population size, with fitted exponential rate in $N$ indistinguishable from zero.
The weak rise from $N = 4$ to $16$ visible in the graze times
($39 \to 72$\,s) does not survive absorption labeling --- it was a graze artifact.
The published-style picture of a lifetime that grows with oscillator count, familiar
from ring-topology chimeras~\cite{wo2011}, is absent here in both label and shape:
adding partials does not buy persistence at this corner, at any tested $N$.

\subsection{The inversion: the high-morph corner never lets go}

The corner the application had classified as unreliable behaves in the opposite way.
At $A = 0.2$, between 51\% and 78\% of seeds (rising with $N$) \emph{never absorb}:
extending the horizon eightfold to $t_{\max} = 16{,}000$\,s leaves the censored
fraction essentially unchanged, so the censoring is irreducible --- a dynamical
finding, not a horizon artifact. The minority that does absorb does so promptly and
increasingly on the first approach: the zero-graze fraction among absorbers rises
from 0.08 at $N = 4$ to 0.96 at $N = 64$. Capture at $A = 0.2$ is bimodal ---
absorb almost immediately, or settle into a persistent, intermittently grazing
chimera~\cite{sudaokuda2015} with no defined absorption time. The stability folklore of the deployment is
therefore inverted: the ``reliable'' corner always dies, and the ``risky'' corner
contains the only immortal chimeras in the system. Section~\ref{sec:reduced} shows
both facts are exactly where the bifurcation diagram says they should be.

\subsection{Aging: hazard rises with every breath}

Absorption at $A = 0.5$ is not memoryless. Censored-Weibull fits give shape
parameters $k_{\mathrm{abs}}$ of $2.27$, $2.22$, $2.49$, $2.67$, $2.44$, $2.72$, and
$3.03$ for the seven sizes $N = 4$--$64$, all confidence intervals well above 1 (Fig.~\ref{fig:weibull}) ---
stronger aging than in the graze times ($k_{\mathrm{graze}} = 0.82$--$2.19$), so the
survival shoulder is not a labeling artifact. The natural per-breath Bernoulli model
--- a constant absorption probability per high-$R$ pass --- is rejected at every
$A = 0.5$ point ($\chi^2$ $p < 0.001$ against geometric cycles-to-absorption;
Weibull-in-cycles shape $k_{\mathrm{cyc}} \approx 2.1$--$2.9$): the per-pass hazard
\emph{rises} with cycle index. The average per-pass absorption probability
$\hat p$ nonetheless tracks the plateau, rising from 0.28 at $N = 4$ to 0.44 by
$N = 24$ and flattening (Fig.~\ref{fig:geometric}), while the mean number of
recovered grazes before death falls from 2.5 to 1.25. At $A = 0.2$ the strong
apparent aging in the graze statistics ($k_{\mathrm{graze}}$ up to 5.3) is retired:
on the few true absorptions $k_{\mathrm{abs}} \approx 0.3 < 1$, a fast, front-loaded
minority rather than aging. The never-absorbers hold a mean incoherent-population
level $\overline{R}_{\mathrm{incoh}} \approx 0.66$--$0.68$, sitting at the stable
fixed point $r^* = 0.678$ identified in Sec.~\ref{sec:reduced}.

\begin{figure}[t]
\centering
\includegraphics[width=0.9\linewidth]{fig3}
\caption{Absorption survival and aging. Kaplan--Meier survival curves by $N$
(left, $A = 0.5$) develop a shoulder with increasing $N$; censored-Weibull shape
$k_{\mathrm{abs}}(N)$ (right) rises from 2.3 to 3.0 with all confidence intervals
above the memoryless value $k = 1$.}
\label{fig:weibull}
\end{figure}

\begin{figure}[t]
\centering
\includegraphics[width=0.85\linewidth]{fig4}
\caption{Per-pass absorption statistics at $A = 0.5$. Mean per-pass absorption
probability $\hat p(N)$ rises then plateaus, mirroring $\tau_{\mathrm{abs}}(N)$;
the geometric (constant-$p$) model is rejected at every point --- hazard rises with
cycle index ($k_{\mathrm{cyc}} > 1$).}
\label{fig:geometric}
\end{figure}

\section{Mechanism: the breath-synchronized ratchet}\label{sec:mechanism}

\subsection{Deaths lock to the breathing phase}

The chimera at $A = 0.5$ breathes: $\min(R_1, R_2)$ oscillates with period
$T_b \approx 21$--$25$\,s, its envelope reaching $\approx 0.83$ against the
$\theta = 0.85$ boundary --- the cycle passes within a hair of near-synchrony once
per breath. Collapse \emph{attempts} (long grazes) lock to this cycle: their phases
relative to the preceding envelope peak cluster at 0.13--0.20 of the cycle, with
resultant length tightening from 0.43 to 0.69 as $N$ grows and pooled Rayleigh
$p \approx 0$. True absorptions inherit the lock: pooled over $A = 0.5$,
$\bar R = 0.21$ with Rayleigh $p = 2.2\times10^{-6}$ (three of six individually
testable points significant; clustering weakens at large $N$, where capture
completes within 2--3 cycles and leaves few phases to test). Mean absorption takes
roughly six breaths ($\tau_{\mathrm{abs}}/T_b \approx 6$). Death, when it comes,
comes at a particular moment of the breath (Fig.~\ref{fig:rose}).

\subsection{The envelope ratchets monotonically to capture}

Within individual runs, successive breathing maxima $M_k$ of $\min(R_1, R_2)$ do not
fluctuate about a level --- they climb. The mean per-cycle increment
$\langle \Delta M \rangle$ is $+0.031, +0.052, +0.078, +0.091$ for
$N = 8, 16, 32, 64$, with cluster-bootstrap confidence intervals strictly positive
at every point; the fraction of runs with monotone (allowing one inversion) maxima
reaches 1.00 by $N = 32$. Where enough sub-threshold approach peaks exist to fit,
$\log(\theta - M_k)$ falls linearly in cycle index with median slopes $-0.55$ to
$-0.77$ and a negative slope in 100\% of fittable runs: the approach to the
synchronization boundary is an exponential spiral-out (Fig.~\ref{fig:ratchet}). The
control rules out artifacts: the same detector and windowing applied to the
$A = 0.2$ never-absorbers over a median of $\sim$200 cycles gives
$\langle \Delta M \rangle \approx -0.001$ and a ratchet fraction of exactly zero ---
stationary, with a slight \emph{relaxation} consistent with settling onto an
attractor.

\subsection{What collapse is not, and what organizes it}

Because the oscillators are strictly identical, a natural mechanism would be escape
transverse to the Poisson (Ott--Antonsen) submanifold~\cite{ott2008,ott2009} on
which the reduced chimera solutions live. We tested this directly with the moment observable
$D = \sum_{m=2}^{4} |Z_m - Z_1^m|^2$, validated in three stages (the Poisson
identity $Z_m = Z_1^m$ to machine precision on M\"obius-constructed states;
dynamical invariance of on-manifold initial conditions under the shipped integrator;
a clustered-state contrast), and the hypothesis fails comprehensively: initial
distance $D_0$ does not predict lifetime ($|\rho| \le 0.15$ at all fourteen
parameter points), event-aligned $\langle D \rangle$ shows no ramp preceding
collapse, and a paired interventional experiment injecting controlled $D_0$ levels
leaves lifetimes flat. Collapse is not transverse escape.

It is instead organized in the collective coordinates. Among static features of the
seed, the initial inter-population phase offset dominates: $|\Delta\phi_0|$
anticorrelates with lifetime at $\rho \approx -0.43$ to $-0.59$ for $N \ge 8$
(median magnitude 0.52), versus the $D_0$ null's $\le 0.08$, and a cross-validated
regression on the $t = 0$ collective features reaches $R^2 \approx 0.2$
(quadratic up to 0.44). Larger initial phase separation between the populations
means earlier death --- a fact Sec.~\ref{sec:reduced} converts from regression to
dynamics.

\begin{figure}[t]
\centering
\includegraphics[width=0.85\linewidth]{fig5}
\caption{Breath-phase locking. Circular histograms of collapse-attempt phases
(left) and true-absorption phases (right) relative to the preceding breathing-envelope
peak, $A = 0.5$. Attempts cluster at 0.13--0.20 of the cycle (pooled Rayleigh
$p \approx 0$); absorptions inherit the lock (pooled $p = 2.2\times10^{-6}$).}
\label{fig:rose}
\end{figure}

\begin{figure}[t]
\centering
\includegraphics[width=0.9\linewidth]{fig6}
\caption{The ratchet. Ensemble-averaged breathing maxima $\langle M_k \rangle$
versus cycle index $k$ at $A = 0.5$ for several $N$, with the $A = 0.2$
never-absorber band (stationary, $\langle \Delta M \rangle \approx -0.001$) for
contrast. Per-cycle increments are strictly positive at every $N$;
$\log(\theta - M_k)$ falls linearly (median slopes $-0.55$ to $-0.77$).}
\label{fig:ratchet}
\end{figure}

\section{Reduced dynamics: the post-homoclinic ghost}\label{sec:reduced}

\subsection{The reduced system, validated against its source}

In the $N \to \infty$ mean-field, the two-population system closes in the
per-population order parameters. Writing $\rho_\sigma e^{-i\phi_\sigma}$ for
population $\sigma$'s centroid, the reduced flow is~\cite{abrams2008}
\begin{equation}
\dot\rho_1 = -\frac{\rho_1^2 - 1}{2}\Big[\mu \rho_1 \cos\alpha
+ \nu \rho_2 \cos(\phi_2 - \phi_1 - \alpha)\Big],
\label{eq:threeD}
\end{equation}
together with the corresponding $\dot\phi_1$ equation and the $1 \leftrightarrow 2$
images; rotational symmetry reduces this to three variables
$(\rho_1, \rho_2, \psi = \phi_1 - \phi_2)$. On the invariant manifold $\rho_1 = 1$
(synchronized population exactly locked) the flow further reduces to two variables
$(r = \rho_2, \psi)$, Eq.~(12) of Ref.~\cite{abrams2008}, whose stationary chimeras,
saddle-node curve $A_{SN}(\beta)$, and Hopf curve $A_H(\beta)$ are known in closed or
series form, with the linear stability of these two-population chimeras
characterized in detail by Laing~\cite{laing_twopop}.

Before applying any of this to our system, the implementation was gated against the
source paper itself: the parametrized stationary chimeras zero the reduced
right-hand side to $4\times10^{-16}$; integration reproduces the published regimes
at $\beta = 0.1$ (stationary at $A = 0.2$; breathing at $A = 0.28$ and $0.35$, with
the longer period at the larger $A$); and numerically located saddle-node and Hopf
points match the published series to $5\times10^{-5}$. Every number below is
parameter-free and, by the integrator's documented clock, in literal seconds.

\subsection{The two corners straddle the bifurcation set}

At $\beta = 0.05$: $A_{SN} = 0.0948$, $A_H = 0.2703$, and the homoclinic --- located
by bisection on the divergence of the breathing cycle's period --- at
$A_{hc} = 0.4096$ (bracket width $8.8\times10^{-5}$). Traced over
$\beta \in [0.03, 0.20]$, the homoclinic curve descends, turns, and rises into the
Takens--Bogdanov point at $(0.2239, 0.3372)$ where it meets the saddle-node and Hopf
curves --- the non-monotone approach is demanded by the termination and serves as an
internal consistency check (Fig.~\ref{fig:diagram}).

The verdicts follow immediately. $A = 0.2$ lies in the stable-chimera band
($A_{SN} < 0.2 < A_H$): the reduced system has a stable stationary chimera at
$r^* = 0.6781$ (integration amplitude decaying to $10^{-9}$), with weak relaxation
$\sigma = -0.0054$\,s$^{-1}$ --- matching the never-absorbers, whose mean
incoherent level sits at $r^*$ and whose envelope drifts slightly \emph{down}, and
explaining the bimodal capture as a basin question: seeds either fall to
synchrony on the first approach or onto the chimera attractor forever. $A = 0.5$
lies beyond the homoclinic ($0.5 > A_{hc} = 0.4096$): \emph{no chimera attractor
exists at the operating corner}. The chimera fixed point is an unstable spiral
($\sigma = +0.01243$\,s$^{-1}$, $\omega = 0.3277$\,s$^{-1}$), the breathing limit
cycle that once surrounded it is destroyed, and sync is globally attracting. What
the deployed voice has been operating is a transient spiraling outward along the
ghost of that destroyed cycle (Fig.~\ref{fig:portraits}).

\begin{figure}[t]
\centering
\includegraphics[width=0.85\linewidth]{fig7}
\caption{Stability diagram of the two-population system in the $(\beta, A)$ plane:
saddle-node (Eq.~17 of Ref.~\cite{abrams2008}) and Hopf (Eq.~18) curves with numeric
checks, the numerically traced homoclinic curve terminating at the Takens--Bogdanov
point $(0.2239, 0.3372)$, and the two operating corners of the deployed system. The
nominally stable corner ($A = 0.5$) is post-homoclinic --- no chimera attractor ---
while the nominally unreliable corner ($A = 0.2$) sits inside the stable-chimera
band.}
\label{fig:diagram}
\end{figure}

\begin{figure}[t]
\centering
\includegraphics[width=0.9\linewidth]{fig8}
\caption{Reduced-system phase portraits at $\beta = 0.05$. Left ($A = 0.2$): stable
chimera spiral at $r^* = 0.678$ with the saddle and its manifolds organizing the
basin. Right ($A = 0.5$): the chimera fixed point is an unstable spiral
($\sigma = +0.0124$\,s$^{-1}$); canonical-seed trajectories spiral out to
synchronization.}
\label{fig:portraits}
\end{figure}

\subsection{Quantitative match with no fitted parameters}

Table~\ref{tab:match} compares the reduced flow to the finite-$N$ measurements. The
spiral's rotation gives a linear breath period $2\pi/\omega = 19.2$\,s and a
nonlinear period of 25.0\,s for trajectories at canonical-seed amplitudes ---
bracketing the measured 21--25\,s. The reduced per-cycle approach slopes,
computed identically to the measurement from each run's seed-mapped initial
condition, are $-0.62, -0.66, -0.66, -0.73$ for $N = 8$--$64$ against measured
$-0.55, -0.50, -0.59, -0.77$: the trend with $N$ is reproduced even though the
reduced flow is $N$-free, because the seed ensembles at different $N$ occupy
different collective initial conditions (mean $R_{\mathrm{incoh},0}$ falls
$\sim 1/\sqrt N$). The single substantial mismatch is the absolute capture time:
the reduced median is 43\,s against the measured 139\,s --- the mean-field flow
dies $3.2\times$ faster than the finite-$N$ system (Sec.~\ref{sec:discussion}).
(A single trajectory started near the focus, by contrast, escapes only at
$\approx 726$\,s: at $\sigma = 0.0124$\,s$^{-1}$ the $\sim$9 e-folds from a
near-focus deviation take exactly that long, while the seed ensemble starts at
finite amplitude --- the two timescales are the same spiral read from different
starting radii.)

\begin{table}[t]
\centering
\caption{Reduced flow versus finite-$N$ measurement at the operating corner
($A = 0.5$, $\beta = 0.05$). No fitted parameters; times in seconds at the
integrator's documented 1:1 clock. The ratio is reduced$/$measured throughout; the
final row reports the $A = 0.2$ chimera level for comparison.}
\label{tab:match}
\begin{tabular}{lccc}
\toprule
quantity & reduced & measured & ratio (reduced/measured) \\
\midrule
breath period $T_b$ & 25.0 (linear 19.2) & 21--25 & $\approx 1.1$ \\
per-cycle spiral slope (pooled) & $-0.67$ & $-0.55$ to $-0.77$ & $\approx 1.0$ \\
median capture time & 43 & 139 & 0.31 \\
$A{=}0.2$: chimera level & $r^* = 0.678$ & $\overline{R}_{\mathrm{incoh}} = 0.66$--$0.68$ & $\approx 1.0$ \\
\bottomrule
\end{tabular}
\end{table}

\subsection{Individual lifetimes from a three-variable model}

Finally, the reduced system predicts \emph{individual} deaths. Mapping every
$A = 0.5$ seed to its collective coordinates
$(\rho_1, \rho_2, \psi)_0 = (R_{\mathrm{sync}}, R_{\mathrm{incoh}},
\Delta\phi)_0$ and integrating the three-variable flow to capture yields a
predicted lifetime whose Spearman correlation with the measured
$t_{\mathrm{abs}}$ is 0.445 pooled (up to 0.60 per $N$) --- comparable to the best
fitted static feature --- and, crucially, carries a \emph{positive partial
correlation of $+0.20$--$0.32$ beyond $|\Delta\phi_0|$ at every $N$}: the
parameter-free dynamics predicts lifetime variance no static regression on the
same inputs captures (Fig.~\ref{fig:scatter}). The transverse coordinate stays
quiet throughout ($\rho_1$ deviating from 1 by $\le 0.003$ median), justifying the
manifold picture a posteriori. This $N$-independence is not merely a mean-field
artifact: for identical oscillators under global sinusoidal coupling,
Watanabe--Strogatz theory~\cite{ws1993,ws1994} reduces each population to a
three-parameter flow at \emph{every} $N$, with the remaining $N-3$ degrees of
freedom frozen as constants of motion --- as the group-theoretic~\cite{mms2009}
and hierarchical-population~\cite{pr2008} formulations make explicit. The
three-variable system integrated here is exact on the uniform-constants (Poisson)
submanifold, where the collective dynamics closes independently of $N$; the
off-manifold constants, which a finite-$N$ ensemble need not realize, are precisely
where the residual $3.2\times$ prolongation must originate --- sharpening, not
softening, the open problem.

\begin{figure}[t]
\centering
\includegraphics[width=0.8\linewidth]{fig9}
\caption{Run-by-run lifetime prediction. Reduced-model capture time from each
seed's collective initial condition versus measured $t_{\mathrm{abs}}$, colored by
$N$ ($A = 0.5$; pooled Spearman $\rho = 0.445$; partial correlation beyond
$|\Delta\phi_0|$ positive at every $N$).}
\label{fig:scatter}
\end{figure}

\section{Discussion}\label{sec:discussion}

\textbf{The inversion, operationally.} The deployed voice's celebrated reliability
--- $\approx 96\%$ in-state over half-hour renders --- is revealed as 100\%
supervisor-maintained: at the operating corner there is nothing to be reliable,
since every unsupervised trajectory dies in $\approx 139$\,s regardless of how many
oscillators are added. Conversely the corner discarded as unreliable holds the
system's only genuine attractor. For applications, two consequences follow. First,
supervision is \emph{intrinsic} to post-homoclinic operation at any $N$ ---
re-seeding is not a finite-size patch but the entire persistence mechanism. Second,
collapse detection should be absorption-grade: a detector with a recovery window
(rather than the 2-s sustain rule, which over-triggers on 69--97\% of firings)
would intervene an order of magnitude less often for the same uptime.

\textbf{The open problem: noise that prolongs.} The reduced flow explains the
geometry of death --- where the corners sit, the breath period, the spiral rate,
even individual lifetimes' ordering --- but captures $3.2\times$ too fast, and the
flat $\tau_{\mathrm{abs}}(N)$ plateau has no analog in an $N$-free collective flow.
Both point the same direction: finite-$N$ fluctuations \emph{extend} the transient's
life relative to mean-field, the opposite of the usual intuition in which noise
destroys metastable states. At the operating corner this prolongation is $N$-independent --- the same factor of
$3.2$ across the full $16\times$ range of $N$ --- though that independence is itself
a property of the corner's depth beyond the homoclinic, not of every post-homoclinic
parameter: at $\beta = 0.10$ it survives at the depth-matched point (coefficient of
variation $0.14$ across $N$) but weakens at the deeper $A = 0.5$, where the factor
falls from $2.0$ to $1.1$ (Appendix~\ref{app:beta}). A natural candidate for the
prolongation is fluctuation-induced re-injection --- finite-size kicks scattering the
trajectory back toward the ghost region the deterministic spiral is leaving. We
tested the simplest version directly (Appendix~\ref{app:noise}): additive Gaussian
noise of amplitude $c/\sqrt N$ on the collective variables of the reduced flow. It
does produce re-injection and prolongation, but the prolongation is
$N$-\emph{dependent} --- strong at small $N$ and absent by $N = 64$ --- because the
noise amplitude is. The observed $N$-independent factor of $3.2$ at the operating
corner therefore rules out additive collective noise and points to a mechanism whose
effective re-injection strength does not scale away as $1/\sqrt N$ (for example
structured or multiplicative fluctuations tied to the breath cycle). We report the
puzzle as measured and leave the responsible mechanism open; the rising-hazard
(aging-in-cycles) structure and the breath-locked phase distribution remain sharp
constraints any such theory must satisfy.

\textbf{Relation to ring-topology transients.} On nonlocally coupled rings,
finite-size chimeras are chaotic transients with lifetimes growing rapidly in
$N$~\cite{wo2011}. The two-population mean-field case is qualitatively different:
its collective flow is $N$-free, the $N$-dependence entering only through the
initial-condition ensemble --- and correspondingly we find no lifetime growth at
all. The contrast suggests the $N$-scaling of chimera transience is a property of
coupling topology rather than of chimeras per se.

\textbf{Scope.} All results are at strictly identical frequencies, equal population
sizes, the shipped phase lag ($\beta = 0.05$, with the phenomenology reproduced at
$\beta = 0.10$; Appendix~\ref{app:beta}), and one canonical seed family; the absolute
$\tau_{\mathrm{abs}}$ carries the documented 22--27\% labeling sensitivity; and the
collective-coordinate summaries are a lossy projection of the $2N$-phase state
(weakest at $N = 4$, where they predict little). Within that scope, every headline
claim is criterion-robust, bit-reproducible, and regenerable from committed
configurations.

\section{Conclusion}\label{sec:conclusion}

An engineering question --- why does a synthesizer's chimera need so much
re-seeding? --- unwound into a dynamical one. The collapse detector turns out to
detect grazes, not deaths: most sustained near-synchrony excursions heal, and true
absorption takes twice as long as first passage suggests. Corrected lifetimes are
$N$-independent and aging, deaths lock to the breathing phase and arrive by a
monotone ratchet of the envelope, and the bifurcation structure of the reduced
two-population model explains why: the operating corner lies beyond the homoclinic,
so the deployed state is a transient on a destroyed cycle's ghost, while the
discarded corner holds the true chimera attractor. A three-variable, parameter-free
flow predicts the period, the spiral rate, and the ordering of individual deaths
--- and underestimates their timing by $3.2\times$, leaving one clean open problem:
at finite $N$, fluctuations make this chimera live longer, not die faster.

\begin{acknowledgments}
This work was conducted independently of the author's institutional appointments.
No external funding supported this work. Generative AI tools (Anthropic Claude and
Claude Code) were used to implement the analysis code and experiment harnesses and
to assist in drafting the manuscript text; all results, derivations, and citations
were verified by the author, who takes full responsibility for the content.
\end{acknowledgments}

\section*{Author Declarations}
\subsection*{Conflict of Interest}
The author has no conflicts to disclose.

\subsection*{Author Contributions}
\textbf{Avery Karlin:} Conceptualization (lead); Methodology (lead); Software
(lead); Formal analysis (lead); Investigation (lead); Data curation (lead);
Visualization (lead); Writing -- original draft (lead); Writing -- review \&
editing (lead).

\section*{Data Availability Statement}
The data and code that support the findings of this study --- the integrator,
experiment harnesses, campaign records (JSONL), committed configurations, and
figure-generation scripts --- are openly available in the annealMusic repository at
github.com/akarlin3/annealMusic and archived at Zenodo
(doi.org/10.5281/zenodo.20564867); python3 tools/paper-figures/run\_all.py
regenerates all figures.

\appendix

\section{Robustness of the phenomenology at $\beta = 0.10$}\label{app:beta}

To show the operating-corner phenomenology is not an artifact of the shipped phase
lag $\beta = 0.05$, we repeated the finite-$N$ campaign at $\beta = 0.10$, the value
at which the reduced model is validated against the published
regimes~\cite{abrams2008}. Reduced-model placement gives a saddle-node at
$A_{SN} = 0.178$, a Hopf at $A_H = 0.278$, and a homoclinic at
$A_{hc}(0.10) = 0.354$ (bisection bracket width $9\times10^{-5}$). We sampled the
post-homoclinic corner $A = 0.5$ (parallel to the operating corner) and, to match its
distance beyond the homoclinic, a depth-matched point
$A = A_{hc} + 0.090 = 0.444$; $A = 0.2$ and a cleaner mid-band point $A = 0.24$ sit
in the stable-chimera band as never-absorber controls.

Every headline result survives (Fig.~\ref{fig:appA}). The absorption lifetime is flat
in $N$ ($\hat\tau_{\mathrm{abs}}(64)/\hat\tau_{\mathrm{abs}}(8) = 0.93$ at $A = 0.5$;
$0.87$ at $A = 0.444$; fitted exponential rate $\approx -1\times10^{-3}$ per
oscillator, zero censoring). The censored-Weibull shape stays above one
($k_{\mathrm{abs}} = 1.8$--$2.8$ at $A = 0.5$, $1.4$--$1.6$ at $A = 0.444$), so hazard
still rises. The per-cycle ratchet is positive
($\langle\Delta M\rangle = +0.06$ to $+0.19$, bootstrap CIs excluding zero), and true
absorptions remain breath-phase-locked (pooled Rayleigh $p = 6.5\times10^{-5}$). The
reduced flow reproduces the breath period ($\approx 23$--$25$\,s, measured
$20$--$25$\,s) and spiral rate. The stable-band controls behave as expected: the
mid-band $A = 0.24$ holds a near-$N$-independent persistent fraction ($44\%$ at
$N = 8$, $52\%$ at $N = 64$; 50 seeds each), a cleaner never-absorber control than
$A = 0.2$, whose persistent fraction swings more widely ($28\%$ to $60\%$) as it sits
close to $A_{SN}$.

One caveat refines a body-text claim. The transient is shallower at $\beta = 0.10$
($\approx 2$ breath cycles to absorption vs $\approx 6$ at the operating corner), and
correspondingly the reduced-to-measured capture ratio --- the $3.2\times$ prolongation
of Sec.~\ref{sec:discussion} --- is $N$-independent only at comparable post-homoclinic
depth. At the depth-matched $A = 0.444$ it stays flat (coefficient of variation $0.14$
across $N$), but at the deeper $A = 0.5$ it falls from $2.0$ to $1.1$ across $N$. The
$N$-independence of the prolongation is thus a property of the operating corner's
depth beyond the homoclinic, not of every post-homoclinic parameter.

\begin{figure}[t]
\centering
\includegraphics[width=0.95\linewidth]{figA}
\caption{Robustness at $\beta = 0.10$. (a)~Absorption lifetime
$\hat\tau_{\mathrm{abs}}(N)$ is flat in $N$ at both the post-homoclinic corner
$A = 0.5$ and the depth-matched $A = 0.444$. (b)~Censored-Weibull shape
$k_{\mathrm{abs}}(N)$ stays above the memoryless value $k = 1$ (aging persists).
(c)~True-absorption breath phases remain clustered (pooled $\bar R = 0.34$, Rayleigh
$p = 6.5\times10^{-5}$, $n = 79$).}
\label{fig:appA}
\end{figure}

\section{A finite-size noise test of fluctuation re-injection}\label{app:noise}

The reduced flow captures $\approx 3.2\times$ faster than the finite-$N$ system, and at
the operating corner that prolongation is $N$-independent
(Sec.~\ref{sec:discussion}). The natural candidate --- fluctuation-induced
re-injection --- predicts that adding finite-size noise to the collective flow should
reproduce both facts. We tested this directly, integrating the three-variable reduced
flow at the operating corner ($A = 0.5$, $\beta = 0.05$) with the Euler--Maruyama
scheme and adding Gaussian noise of amplitude $\sigma_{\mathrm{noise}} = c/\sqrt N$ to
the collective variables $(\rho_1, \rho_2, \psi)$, with reflection keeping $\rho$ in
$[0,1]$. At $c = 0$ the scheme reproduces the deterministic capture times to a median
$0.25\%$ per initial condition. The physical scale, estimated from the measured
high-frequency fluctuation of $R_{\mathrm{incoh}}$ in the baseline data, is
$c \approx 0.05$. We swept $N \in \{8,16,32,64\}$, $c \in \{0, 0.025, 0.05, 0.1,
0.2\}$, $200$ realizations each, from the Sec.~\ref{sec:reduced} seed-mapped initial
conditions.

The result is a clean negative (Fig.~\ref{fig:appB}). At the physical scale
$c \approx 0.05$ the noise barely perturbs capture (prolongation factor
$\approx 0.97$, $N$-independent). Larger amplitudes do prolong --- at $c = 0.2$ the
trajectory is repeatedly re-injected, accumulating up to ${\sim}240$ breath cycles
before capture --- but the prolongation is strongly $N$-\emph{dependent} (factor
$2.4$ at $N = 8$ falling to $0.7$ at $N = 64$; coefficient of variation $0.40$ across
$N$), and never reaches the measured $\approx 3.2$ at any cell where it is
$N$-independent. The rising-hazard constraint holds only at small $c$ ($k > 1$ for
$c \le 0.05$, $k \approx 1$ for $c \ge 0.1$), while breath-phase locking survives
throughout. Thus additive $1/\sqrt N$ noise on the collective flow reproduces the
mechanism (re-injection) but with the wrong $N$-scaling: because its amplitude is
$N$-dependent, so is the prolongation it produces. The $N$-independent $3.2\times$
factor requires a different mechanism.

\begin{figure}[t]
\centering
\includegraphics[width=0.85\linewidth]{figB}
\caption{Finite-size noise test. Prolongation factor (median capture relative to the
deterministic reduced flow) versus additive-noise amplitude $c$
($\sigma_{\mathrm{noise}} = c/\sqrt N$) at $A = 0.5$, $\beta = 0.05$. At the physical
scale $c \approx 0.05$ the curves sit near unity; at larger $c$ they fan out with $N$,
so the prolongation is $N$-dependent and never reaches the measured $\approx 3.2$
(dotted) where it is $N$-independent.}
\label{fig:appB}
\end{figure}

\begin{thebibliography}{99}
\bibitem{kb2002} Y.~Kuramoto and D.~Battogtokh, ``Coexistence of coherence and incoherence in nonlocally coupled phase oscillators,'' Nonlinear Phenom.\ Complex Syst.\ \textbf{5}, 380 (2002).
\bibitem{as2004} D.~M.~Abrams and S.~H.~Strogatz, ``Chimera states for coupled oscillators,'' Phys.\ Rev.\ Lett.\ \textbf{93}, 174102 (2004).
\bibitem{abrams2008} D.~M.~Abrams, R.~Mirollo, S.~H.~Strogatz, and D.~A.~Wiley, ``Solvable model for chimera states of coupled oscillators,'' Phys.\ Rev.\ Lett.\ \textbf{101}, 084103 (2008); erratum \textbf{101}, 129902 (2008).
\bibitem{wo2011} M.~Wolfrum and O.~E.~Omel'chenko, ``Chimera states are chaotic transients,'' Phys.\ Rev.\ E \textbf{84}, 015201(R) (2011).
\bibitem{pa2015} M.~J.~Panaggio and D.~M.~Abrams, ``Chimera states: coexistence of coherence and incoherence in networks of coupled oscillators,'' Nonlinearity \textbf{28}, R67 (2015).
\bibitem{om2018} O.~E.~Omel'chenko, ``The mathematics behind chimera states,'' Nonlinearity \textbf{31}, R121 (2018).
\bibitem{panaggio2016} M.~J.~Panaggio, D.~M.~Abrams, P.~Ashwin, and C.~R.~Laing, ``Chimera states in networks of phase oscillators: The case of two small populations,'' Phys.\ Rev.\ E \textbf{93}, 012218 (2016).
\bibitem{lem2019} N.~Lem and Y.~Orlarey, ``Kuroscillator: a Max-MSP object for sound synthesis using coupled-oscillator networks,'' in \emph{Proceedings of the 14th International Symposium on Computer Music Multidisciplinary Research (CMMR)} (Marseille, France, 2019).
\bibitem{lemsmc} N.~Lem, ``Sound in multiples: synchrony and interaction design using coupled-oscillator networks,'' in \emph{Proceedings of the 16th Sound and Music Computing Conference (SMC)} (M\'alaga, Spain, 2019), pp.\ 470--475.
\bibitem{kuramoto1984} Y.~Kuramoto, \emph{Chemical Oscillations, Waves, and Turbulence} (Springer-Verlag, Berlin, 1984).
\bibitem{sk1986} H.~Sakaguchi and Y.~Kuramoto, ``A soluble active rotator model showing phase transitions via mutual entrainment,'' Prog.\ Theor.\ Phys.\ \textbf{76}, 576 (1986).
\bibitem{sudaokuda2015} Y.~Suda and K.~Okuda, ``Persistent chimera states in nonlocally coupled phase oscillators,'' Phys.\ Rev.\ E \textbf{92}, 060901(R) (2015).
\bibitem{ott2008} E.~Ott and T.~M.~Antonsen, ``Low dimensional behavior of large systems of globally coupled oscillators,'' Chaos \textbf{18}, 037113 (2008).
\bibitem{ott2009} E.~Ott and T.~M.~Antonsen, ``Long time evolution of phase oscillator systems,'' Chaos \textbf{19}, 023117 (2009).
\bibitem{laing_twopop} C.~R.~Laing, ``Dynamics and stability of chimera states in two coupled populations of oscillators,'' Phys.\ Rev.\ E \textbf{100}, 042211 (2019).
\bibitem{pr2008} A.~Pikovsky and M.~Rosenblum, ``Partially integrable dynamics of hierarchical populations of coupled oscillators,'' Phys.\ Rev.\ Lett.\ \textbf{101}, 264103 (2008).
\bibitem{mms2009} S.~A.~Marvel, R.~E.~Mirollo, and S.~H.~Strogatz, ``Identical phase oscillators with global sinusoidal coupling evolve by M\"obius group action,'' Chaos \textbf{19}, 043104 (2009).
\bibitem{ws1993} S.~Watanabe and S.~H.~Strogatz, ``Integrability of a globally coupled oscillator array,'' Phys.\ Rev.\ Lett.\ \textbf{70}, 2391 (1993).
\bibitem{ws1994} S.~Watanabe and S.~H.~Strogatz, ``Constants of motion for superconducting Josephson arrays,'' Physica D \textbf{74}, 197 (1994).
\end{thebibliography}

\end{document}
